% =====================================================================
%  APÊNDICE B — código dos modelos (27/10/2026, pedido do autor depois da
%  leitura da orientadora: "a classe arima, classe garch, selecao e
%  estrategia", "da forma mais simples e crua possível, sem logs e sem
%  flags de teste adicional").
%  Os arquivos listados são gerados pelo Makefile a partir de robos/Include/*.mqh
%  por ferramentas/scripts/codigo_apendice.sh, com os mesmos marcadores do
%  EA: saem comentários, mensagens de erro, a montagem do relatório da
%  seleção e as opções experimentais (t de Student, BIC, máxima
%  verossimilhança, janela própria do GARCH). A densidade normal, que mora
%  em GarchDistribuicao.mqh, entra no Garch.mqh por //@apendice.
%  O EA reduzido mais os quatro arquivos reduzidos compilam juntos
%  (0 erros, 0 avisos, 27/10/2026).
% =====================================================================
\chapter{Código dos modelos}

Este apêndice reúne os arquivos chamados pelo \textit{Expert Advisor} do
Apêndice~A que contêm os modelos. Como no Apêndice~A, foram retirados os
comentários, as mensagens de registro e as opções experimentais que não fazem
parte deste trabalho. As funções auxiliares de mínimos quadrados e dos testes
estatísticos, assim como as rotinas de envio de ordens, ficam em arquivos
próprios, que não são reproduzidos aqui.

\section{Classe do ARIMA}

A classe \texttt{CArima} diferencia a série, estima os coeficientes pelo método
de Hannan e Rissanen, calcula os resíduos e o AIC e faz a previsão do período
seguinte.

\lstinputlisting[style=mql5]{codigo/Arima.mqh}

\section{Seleção das ordens}

A função \texttt{DeterminarOrdemD} aplica o teste ADF em sequência para definir
a ordem $d$, e a função \texttt{BuscarMelhorPQ} ajusta todas as combinações de
$p$ e $q$ de 0 a 4 e fica com a de menor AIC.

\lstinputlisting[style=mql5]{codigo/SelecaoModelo.mqh}

\section{Classe do GARCH}

A classe \texttt{CGarch} estima os parâmetros do GARCH por máxima
verossimilhança, com as restrições passadas ao otimizador da biblioteca ALGLIB,
e prevê a volatilidade do período seguinte.

\lstinputlisting[style=mql5]{codigo/Garch.mqh}

\section{Decisão de operação}

A função \texttt{DecidirOperacao} recebe a previsão do ARIMA, que dá a direção,
e a volatilidade prevista pelo GARCH, que pode vetar a operação no filtro ou
mudar o tamanho da posição no ajuste.

\lstinputlisting[style=mql5]{codigo/Estrategia.mqh}
